\documentclass{ltjsbook}
\usepackage{docmute} 
\usepackage{../../myStyle} 

\begin{document}
\chapter{おわりに}
\label{H11_Epilogue}
この講義では，心理学における測定，データの扱いについて論じてきました。以下に各章の流れをまとめてみたいと思います。

第\ref{H01_PsychlogicalScaling}章では心理学が何を研究対象にしているのかを考え，態度尺度と呼ばれる方法で作られる数値はなぜそれで間隔尺度水準の数値になるといえるのか，その理論的裏付けについて考えました。
第\ref{H02_TestTheory}章では心理学と少し違った角度，テスト理論という文脈で，目に見えないものを測定するモデルとしてどのようなものが考えられてきたか，古典的テスト理論と現代テスト理論に言及しながら考えてきました。
第\ref{H03_FA}章では一転して心理学に戻り，心理学で最もよく使われている測定モデルである因子分析法について，そのモデルと実践例について考えてきました。
第\ref{H04_Algebra},\ref{H05_MatrixForm}章は線形代数について学びました。数学的側面なので苦手意識を克服できなかった人もいるかも知れませんが，ここで正方行列のスペクトル分解を理解しておくと，その後のデータの扱い一般についての理解が格段に広がるため，どこかで触れなければならなかったのです。上手く理解できなければ，こちらの伝え方の方に問題があります。自分を責めずに，わかるまでトライしてもらえれば幸いです。
第\ref{H06_Attitude}章ではここまでの知識を踏まえた上で，実際に心理学の研究において心理尺度がどのように使われているか，言語による測定がどのようなものかを見てみました。態度尺度のような理論的裏付けが及ばない領域での利用が多くみられ，翻ってそれぞれの領域で考えるべき点，守られるべき点や，どのような改善方法が考えられるかについて検討しました。
第\ref{H07_Distance}章では，心理尺度に至る前のプリミティブな心の判断，類似性をデータにした場合にどのような分析が可能かについて，クラスター分析や多次元尺度構成法について学びました。
第\ref{H08_MDS}章では多次元尺度構成法の中でも，非計量的多次元尺度法や個人差モデル，展開法，非対称多次元尺度法などの応用例を見てきました。これらの方法は心理尺度による研究よりもデータが頑健で，個人差についての明示的なモデルがあり，尺度に対する反応も異なる観点から捉えてモデル化できること，対人関係のような非対称性が重要な側面でも応用可能なモデルがあることなど，その広がりについて論じました。
第\ref{H09_dual}章では双対尺度法を中心に，名義尺度水準のデータを時限解析する方法についてみていきました。ここに至って「カテゴリに数値を付与する」という明確な方向性を改めて確認し，分析者にとって意味のある数値化についてのいくつかの応用モデルもみていきました。
第\ref{H10_Bayesian}章ではモデリングアプローチから，そもそもデータがどのような機序で生成されているか，データ生成メカニズムの観点から分析する方法を見てきました。数理モデルをデータと照合するための便利な方法として，ベイズ統計の枠組みである確率的プログラミング言語の助けを借りながら，研究関心を直接表現することでデータドリブンの研究スタイルを転換させる手段を導入しました。またこの手法から従来法を眺め直すと，いろいろ不便なところも見えてきましたが，両者の良いところを見ながら使うことで今後の分析が楽しくなれば良いなと思っています。

ここでお話しした内容は，「心理学統計法」という枠組みというより「心理測定法」の枠組みに入る内容です。測定法といえば，心理学研究法の下位分野と考えている人もいるかもしれません。

心理学界隈では昨今，公認心理師という国家資格がつくられたことにより，資格対応を考える大学ではカリキュラムが標準化される動きがありました。心理学は非常に幅広い領域をカバーする学問ですし，これに医療や法律の観点を加えて4年間で履修するプログラムを考えるとなると，各領域をどのようにまとめるか，その境界を引いたり複合領域をどう扱うかという問題に直面することになります。当然カリキュラムを考えるところでは激しい討論が重ねられたでしょうし，それによって提供された現状のものが完璧であるとはいえないと思います。しかしそれは資格取得を目指す学生側には関係のないことで，学ぶ人たちはそれぞれ自らの思いを持って，誠実に決められたことを学ぶしかないでしょう。

残念ながら，心理学の各領域の専門家と言われる人たちも，その領域のあらゆるジャンルをマスターしているわけではありません。教える側にも単元ごとに得手不得手があり，よく知っているところはうまく教えられるでしょうが，よく知らないところは教科書の表面的なところをなぞるだけ，ということになるかもしれません。心理統計の領域はそもそも専門家の数が少ないのですが，それに比べてユーザの数，手練れの経験者の数は非常に多いので，ややもするとノウハウ寄りの教え方をする人が多いかも知れません。また教わる側も，心理学と数学を同時に愛する人は少なく，そもそも数学嫌いだからこそ文系たる心理学に行く，という人もいるでしょうから，ややもするとノウハウのみを身につけて，試験をクリアすることを目的にする人が出てくるでしょう。測定の最初に学んだように，資格やその試験など数値目標を設定すると，それを獲得するためのルートは必ずハックされ，チートを探す人が出てきます。公認心理師の試験は膨大な領域に及びますから，合格するための良い戦略は「心理統計の問題が出たら捨てる」かもしれません。

これは「資格」というもの全体の問題で，資格になるからこそ専門性が下がるというパラドキシカルな現象はどこでもあり得る話です。特に心理統計に限って声高にいうべきことではないのかもしれません。
しかしここであえていいたいのは，\textbf{心理測定は心理統計や心理学研究法の下位領域ではなく，心理学そのものである}と筆者は考えるからです。ノウハウに振り切って教えられることで，原理原則が無視されたり軽んじられ，そもそもの発想からは\kenten{合理的な裏付けがえられないほどに}奇妙な研究実践が蔓延してしまっている様を，皆さんも垣間見たのではないでしょうか。資格や科目の一部に矮小化されたり，優先順位が下がることは仕方がないこととはいえ，\textbf{教えられなかったものは存在しなくなる}可能性があります。古い心理学の全集などをみると，かならず数理心理学，計量心理学で一冊は書かれるほどリッチな研究と考察が重ねられていました。今や実験実習の一つの単元で尺度作りが行われるほど，そのノウハウは一般化しましたが，その数字が意味するものが何なのか，誰も知らなくなる日がそこまできているのかもしれません。

心理学が心の学問である，行動の学問であるというのはジョークの一種であり，実は空想の世界のSF小説で，架空の世界で架空の数値を使って，さらに架空の統計量を使って面白おかしくストーリーを物語っているだけだ，というのであれば，それはそれで構わないと筆者は思っています。心理学以外の領域，特にハードサイエンスの人から見たら，私たちがやっていることは総じてコントなのかもしれません。いや違う，心を，人間を，行動を測って理解しようとしているのだというのであれば，馬鹿馬鹿しくも蔓延したルーチンワークのような尺度作成は捨て去って，より建設的な，人の心の深淵に触れうるための学問を積み重ねていかなければならないと思いませんか。

大事なものがこぼれ落ちていく世の中で，私のささやかな足掻きが読者のみなさんにちょっとした引っかかりを生むことができれば，望外の喜びです。

\vspace{2cm}
このテキストは，広島大学総合科学部に非常勤講師としてお招きいただき，これまで筆者の学び，考えてきたことを「測定法」という観点からまとめ直したものになります。他の資料の流用・併用も多くありますが，新たに書き下ろしたセクションも少なくありません。類書として，筆者の現所属での授業資料である「心理学データ解析基礎」，「心理学データ解析応用」のテキストや，数学ワークショップのためにまとめた「心理学者のための線形代数」などがあります。いずれも心理統計教育教材のサイト\footnote{\url{https://kosugitti.github.io/psychometrtics_syllabus/}}にてPDF他公開しておりますので，よろしければそちらもご参照ください。これらのテキストを紙媒体でまとめたものも，実費で販売しております。


\end{document}